% ─────────────────────────────────────────────────────────────────────────────
\subsection*{Overview}

The tick-counter model simulates a periodic tick source that is counted
over time and periodically registered and reset.  Two generators
produce events at different rates: a fast generator fires every
$1/10$~sec (the \emph{tick}) and a slow generator fires every $1$~sec
(the \emph{reset}).  A counter model accumulates incoming ticks and,
upon receiving a reset signal, outputs the current count and resets to
zero.  Ideally the counter outputs $10$ at every reset.

Under PDEVS, the tick and reset generators are simultaneously imminent at
every integer second.  Cadmium uses separate typed ports for the two
signals, enforcing the routing contract at compile time rather than by
value encoding.  The experiment reveals whether the chosen TIME type
preserves the simultaneity invariant.

% ─────────────────────────────────────────────────────────────────────────────
\subsection*{Model Diagram}

\begin{center}
\begin{tikzpicture}[
  box/.style={draw, rectangle, minimum width=1.8cm, minimum height=0.9cm, align=center},
  >=Stealth
]
  \node[box] (g10) at (0, 0.9) {$G_{1/10}$};
  \node[box] (g1)  at (0,-0.9) {$G_1$};
  \node[box] (k)   at (4.5, 0) {$K$};

  \draw[->] (g10) -- node[above, font=\small] {\texttt{add}: \texttt{int} 1} (k);
  \draw[->] (g1)  -- node[below, font=\small] {\texttt{reset}: \texttt{reset\_tick}} (k);
  \draw[->] (k)   -- node[above, font=\small] {\texttt{sum}: \texttt{int} $n$}
                     (6.5, 0) node[right, font=\small] {\texttt{out}};
\end{tikzpicture}
\end{center}

% ─────────────────────────────────────────────────────────────────────────────
\subsection*{Coupled Model}

\[
  C = \langle X_C,\, Y_C,\, D,\, \{M_d\},\, \text{EIC},\, \text{EOC},\, \text{IC} \rangle
\]
\[
  D = \{G_{1/10},\; G_{1},\; K\},\quad X_C = \emptyset,\quad Y_C \subseteq \mathbb{N}
\]
\begin{align*}
  \text{EIC}  &= \emptyset \\
  \text{IC}   &= \{(G_{1/10}.\texttt{out} \to K.\texttt{add}),\;
                   (G_{1}.\texttt{out} \to K.\texttt{reset})\} \\
  \text{EOC}  &= \{(K.\texttt{sum} \to C.\texttt{out})\}
\end{align*}

Expected invariant: every output of $K$ equals $10$.

% ─────────────────────────────────────────────────────────────────────────────
\subsubsection*{Atomic Model Definitions}

\paragraph{Generator $G_{1/10}$.}

\[
  G_{1/10} = \langle X,\; Y,\; S,\;
               \delta_{\text{ext}},\; \delta_{\text{int}},\; \delta_{\text{con}},\;
               \lambda,\; ta \rangle
\]
\begin{align*}
  X &= \emptyset \\
  Y_{\texttt{out}} &= \{1\} \\
  S &= \bigl[0,\;\tfrac{1}{10}\bigr) \\
  \delta_{\text{ext}} &= \bot
    \quad\text{(undefined; } X = \emptyset\text{)} \\
  \delta_{\text{int}}(s) &= 0 \\
  \delta_{\text{con}} &= \bot
    \quad\text{(undefined; } X = \emptyset\text{)} \\
  \lambda(s) &= (\texttt{out}: 1) \\
  ta(s) &= \bigl(\tfrac{1}{10} - s\bigr)\,\text{sec}
\end{align*}

\noindent\textit{Implementation.}
\texttt{basic\_model::pdevs::generator<TIME>}
from namespace \texttt{cadmium}; see \texttt{tick\_gen.hpp}.
$S$ represents elapsed time; since $X = \emptyset$, $s = 0$ always
holds at transition time and $ta$ returns the constant
$\tfrac{1}{10}\,\text{sec}$.
Port \texttt{out}: \texttt{out\_port<int>}; value transported is $1$ (tick).
No EIC connection --- $X = \emptyset$.

\paragraph{Generator $G_1$.}

\[
  G_1 = \langle X,\; Y,\; S,\;
          \delta_{\text{ext}},\; \delta_{\text{int}},\; \delta_{\text{con}},\;
          \lambda,\; ta \rangle
\]
\begin{align*}
  X &= \emptyset \\
  Y_{\texttt{out}} &= \{\texttt{reset\_tick}\} \\
  S &= \bigl[0,\; 1\bigr) \\
  \delta_{\text{ext}} &= \bot
    \quad\text{(undefined; } X = \emptyset\text{)} \\
  \delta_{\text{int}}(s) &= 0 \\
  \delta_{\text{con}} &= \bot
    \quad\text{(undefined; } X = \emptyset\text{)} \\
  \lambda(s) &= (\texttt{out}: \texttt{reset\_tick}) \\
  ta(s) &= \bigl(1 - s\bigr)\,\text{sec}
\end{align*}

\noindent\textit{Implementation.}
\texttt{basic\_model::pdevs::generator<TIME>}
from namespace \texttt{cadmium}; see \texttt{reset\_gen.hpp}.
$S$ represents elapsed time; $ta$ returns the constant $1\,\text{sec}$.
Port \texttt{out}: \texttt{out\_port<reset\_tick>}; value is a
\texttt{reset\_tick} sentinel (unit type with no payload --- the signal
is the arrival itself).
No EIC connection --- $X = \emptyset$.

\paragraph{Counter $K$.}

\[
  K = \langle X,\; Y,\; S,\;
       \delta_{\text{ext}},\; \delta_{\text{int}},\; \delta_{\text{con}},\;
       \lambda,\; ta \rangle
\]
\begin{align*}
  X_{\texttt{add}}   &= \mathbb{N}
    & X_{\texttt{reset}} &= \{\texttt{reset\_tick}\} \\
  Y_{\texttt{sum}}   &= \mathbb{N} \\
  S &= \mathbb{N} \times \{\text{idle},\; \text{ready}\}
\end{align*}

Let $x_{\texttt{add}}$ denote the PDEVS bag on port add,
$x_{\texttt{reset}}$ the bag on port reset.

\begin{align*}
  ta(\bullet,\text{idle})  &= +\infty\,\text{sec}
  & ta(\bullet,\text{ready}) &= 0\,\text{sec}
\end{align*}
\[
  \delta_{\text{ext}}((n, \text{mode}), e, x) =
    \begin{cases}
      (n + |x_{\texttt{add}}|,\; \text{ready}) & x_{\texttt{reset}} \neq \emptyset \\
      (n + |x_{\texttt{add}}|,\; \text{idle})  & x_{\texttt{reset}} = \emptyset
    \end{cases}
\]
\begin{align*}
  \delta_{\text{int}}(n, \text{ready}) &= (0,\; \text{idle}) \\[4pt]
  \delta_{\text{con}}(s, e, x) &=
    \delta_{\text{ext}}\!\left(\delta_{\text{int}}(s),\; 0,\; x\right) \\[4pt]
  \lambda(n, \text{ready}) &= (\texttt{sum}: n)
\end{align*}

\noindent\textit{Implementation.}
\texttt{basic\_model::pdevs::accumulator<int,TIME>}
from namespace \texttt{cadmium}.
$S$ is stored as \texttt{(n:~int,~mode)} with an internal enum
\texttt{\{idle,~ready\}}.
Port \texttt{add}: \texttt{in\_port<int>}; each arriving \texttt{int}
increments $n$.
Port \texttt{reset}: \texttt{in\_port<reset\_tick>}; arrival triggers
the ready state regardless of payload.
Port \texttt{sum}: \texttt{out\_port<int>}; emits $n$ on internal
transition.

% ─────────────────────────────────────────────────────────────────────────────
\subsection*{Simulation Setup}

Both TIME variants are driven by the Cadmium \texttt{pdevs::runner}
with a fixed target time limit.  The runner advances the simulation
until logical time reaches the limit.  If the two runs produce logs of
different length (e.g.\ due to causality breaks in the float run), the
longer log is truncated to the shorter run for comparison.

\paragraph{TIME types exercised.} \texttt{float} and a rational type.

% ─────────────────────────────────────────────────────────────────────────────
\subsection*{Expected Observations}

\textbf{Float time:} IEEE~754 cannot represent $1/10$ exactly.
Floating-point error shifts occasional $G_{1/10}$ ticks across a
1-second boundary; $K$ outputs $9$ or $11$.  Errors may break the
causality chain and are unbounded.
Observed: ${\approx}11.4\%$ error rate over 100\,000 resets
(histogram: $9 \times 10\,406$, $10 \times 88\,624$, $11 \times 970$).

\textbf{Rational time:} $G_{1/10}$ and $G_1$ are always simultaneously
imminent at integer seconds; confluence correctly resolves the tie and
$K$ always outputs $10$.
